\chapter{Mack's published numerical example}

\begin{definition}[Taylor-Ashe cumulative triangle]\label{def:taylorAshe}\lean{VerifiedReserving.taylorAshe}\leanok
The 55 observed cumulative claims in Mack's Table 1 are represented exactly as
integers. Unobserved cells are zero and do not enter any estimator.
\end{definition}

\begin{theorem}[Published development factors]\label{thm:taylorAsheFhat}\lean{VerifiedReserving.taylorAshe_fhat_0, VerifiedReserving.taylorAshe_fhat_1, VerifiedReserving.taylorAshe_fhat_2, VerifiedReserving.taylorAshe_fhat_3, VerifiedReserving.taylorAshe_fhat_4, VerifiedReserving.taylorAshe_fhat_5, VerifiedReserving.taylorAshe_fhat_6, VerifiedReserving.taylorAshe_fhat_7, VerifiedReserving.taylorAshe_fhat_8}\leanok
\uses{def:taylorAshe, def:fhat}
The nine development factors computed from Table 1 equal exact rational
numbers whose decimal roundings are the values printed by \cite{Mack1993}.
\end{theorem}
\begin{proof}\leanok Each equality follows by normalization of the finite sums and division. \end{proof}

\begin{theorem}[Published estimable variances]\label{thm:taylorAsheSigma2}\lean{VerifiedReserving.taylorAshe_sigma2_0, VerifiedReserving.taylorAshe_sigma2_1, VerifiedReserving.taylorAshe_sigma2_2, VerifiedReserving.taylorAshe_sigma2_3, VerifiedReserving.taylorAshe_sigma2_4, VerifiedReserving.taylorAshe_sigma2_5, VerifiedReserving.taylorAshe_sigma2_6, VerifiedReserving.taylorAshe_sigma2_7}\leanok
\uses{def:taylorAshe, def:sigma2, thm:taylorAsheFhat}
The eight variance values supplied by Mack's variance estimator equal exact
rationals with the printed decimal roundings. The ninth value is specified by
an extrapolation convention and is not asserted here as an estimator value.
\end{theorem}
\begin{proof}\leanok Each equality follows by normalization of the weighted squared deviations. \end{proof}

\begin{definition}[Total reserve for Table 1]\label{def:taylorAsheTotalReserve}\lean{VerifiedReserving.taylorAsheTotalReserve}\leanok
\uses{def:taylorAshe, def:reserve}
The total reserve is the sum of the ten chain-ladder reserves computed from
the published triangle.
\end{definition}

\begin{theorem}[Published total reserve]\label{thm:taylorAsheReserve}\lean{VerifiedReserving.taylorAshe_total_reserve_rounds}\leanok
\uses{def:taylorAsheTotalReserve, thm:taylorAsheFhat}
The exact total reserve is strictly between 18,680.5 and 18,681.5 thousand,
so it rounds to the 18,681 thousand printed in Mack's Table 2.
\end{theorem}
\begin{proof}\leanok Normalization of the ten reserve terms proves both strict inequalities. \end{proof}
